% ============ 1 引言 (原书 pp.7-15) ============
\chapter{引言}\label{sec:1}
\index{平展上同调}
\index{Betti 数}
\index{代数拓扑}
\index{层上同调}
\index{Zariski 拓扑}
\index{平展态射}
\index{比较定理}

对复数域上的簇 $X$，$X(\mathbb{C})$ 从 $\mathbb{C}$ 获得一个拓扑，因此可以把代数拓扑的机器用于它的研究.例如，可以定义 $X$ 的 Betti 数 $\beta_r(X)$ 为向量空间 $H^{r}(X(\mathbb{C}), \mathbb{Q})$ 的维数，并且诸如 Lefschetz 不动点公式之类的定理都是可用的.

对任意代数闭域 $k$ 上的簇 $X$，只有 Zariski 拓扑，它太粗糙(即开子集太少)，以至于代数拓扑的方法无法使用.例如，若 $X$ 不可约，则用 Zariski 拓扑计算的群 $H^{r}(X, \mathbb{Z})$ 对一切 $r > 0$ 都为零.

在 1940 年代，Weil 观察到，他在有限域上某些簇(曲线，阿贝尔簇，对角超曲面等)的点数方面的某些结果，可以用一种上同调理论的存在性来解释：该理论给出特征零域上的向量空间，并且对这种理论 Lefschetz 不动点公式成立.他的结果预测了 $\mathbb{C}$ 上 $\mathbb{P}^{d+1}$ 中对角超曲面的 Betti 数公式，后来被 Dolbeault 验证.

大约在 1958 年，Grothendieck 定义了概形的平展``拓扑''，平展上同调理论由他在 M. Artin 和 J.-L. Verdier 的协助下建立.整个理论紧密模仿拓扑空间上通常的层论及其导出函子上同调.对 $\mathbb{C}$ 上的簇 $X$，当 $\Lambda$ 是有限群，$\ell$ 进整数环 $\mathbb{Z}_{\ell}$ 或 $\ell$ 进数域 $\mathbb{Q}_{\ell}$ 时(但对 $\Lambda = \mathbb{Z}$ 不成立)，平展上同调群 $H^{r}(X_{\mathrm{et}}, \Lambda)$ 与复群 $H^{r}(X(\mathbb{C}), \Lambda)$ 一致.当 $X$ 是域 $K$ 的谱时，$X$ 的平展上同调理论恰与 $K$ 的 Galois 上同调理论一致.因此，平展上同调弥合了第一种情形(纯几何的)与第二种情形(纯算术的)之间的鸿沟.

正如我们将在课程中看到的，平展上同调确实给出预期的 Betti 数.此外，它满足 Eilenberg-Steenrod 公理的类比，Poincar\'e 对偶定理，Lefschetz 不动点公式，Leray 谱序列等等.Goresky 和 MacPherson 的相交上同调有平展类比，它为奇异簇提供 Poincar\'e 对偶定理.平展上同调在解释 Weil 的观察方面取得了辉煌的成功.

\subsection*{代数拓扑}

我们简要回顾平展上同调所模仿的理论的起源.

代数拓扑起源于 19 世纪末 Riemann,Betti 和 Poincar\'e 关于``同调数''的工作.在 Emmy Noether 的一个观察之后，焦点转向``同调群''.到 1950 年代，已有若干种给拓扑空间配备(上)同调群的方法，例如 Veblen,Alexander 和 Lefschetz 的奇异同调群，Lefschetz 的相对同调群，Vietoris 同调群，\v{C}ech 同调群，以及 Alexander 上同调群.

Eilenberg 和 Steenrod 1953\footnote{\emph{Foundations of Algebraic Topology}, Princeton.}极大地澄清了这种局面，他们证明：对任意``容许的''拓扑空间对范畴，恰好存在一种上同调理论，满足一组简短的公理.考虑例如这样的范畴：对象是空间对 $(X, Z)$，其中 $X$ 是局部紧拓扑空间，$Z$ 是 $X$ 的闭子集，态射是空间对的连续映射.此范畴上的上同调理论是一个反变函子，它给每个空间对配上一列 Abel 群与映射
\[
\cdots \to H^{r-1}(U) \to H^{r}_{Z}(X) \to H^{r}(X) \to H^{r}(U) \to \cdots, \qquad U = X \setminus Z,
\]
满足下列公理：
\begin{enumerate}
  \item (正合性公理) 上述序列正合；
  \item (同伦公理) 映射 $f^{*}$ 只依赖于 $f$ 的同伦类；
  \item (切除) 若 $V$ 在 $X$ 中开且其闭包与 $Z$ 不相交，则包含映射 $(X \setminus V, Z) \to (X, Z)$ 诱导同构 $H^{r}_{Z}(X) \to H^{r}_{Z}(X \setminus V)$；
  \item (维数公理) 若 $X$ 由单点构成，则 $H^{r}(P) = 0$ ($r \neq 0$).
\end{enumerate}
拓扑学家通常把群 $H^{r}_{Z}(X)$ 写为 $H^{r}(X, U)$.同调理论的公理与上述类似，只是所有箭头的方向都反转.若 $(X, Z) \mapsto H^{r}_{Z}(X)$ 是上同调理论，使得 $H^{r}_{Z}(X)$ 是局部紧 Abel 群(例如离散的或紧的)，则 $(X, Z) \mapsto H^{r}_{Z}(X)^{\vee}$ (Pontryagin 对偶)是同调理论.在这种方法中，隐含地只有一个系数群.

在 1940 年代，Leray 试图理解两个空间 $X$ 与 $Y$ 的上同调群之间的关系，其间给定了连续映射 $Y \to X$.这引导他引入层(系数群的局部系统)，层上同调和谱序列(几乎同时 Roger Lyndon 也独立地做着类似的事，他试图理解群 $G$，正规子群 $N$ 与商群 $G/N$ 的上同调之间的关系).

导出函子在 Cartan 和 Eilenberg 1956\footnote{\emph{Homological Algebra}, Princeton.}中被系统地使用，而 Eilenberg 的学生 Buchsbaum 在 1955 年的博士论文中定义了 Abel 范畴的概念，并把 Cartan-Eilenberg 的导出函子理论推广到这种范畴.(``abelian category'' 这个名称归功于 Grothendieck.)

最后，Grothendieck 在他著名的 1957 年 T\^ohoku 论文\footnote{\emph{Sur quelques points d'alg\`ebre homologique}, T\^ohoku Math. J. 9, 119--221.}中证明，拓扑空间上的 Abel 群层范畴是具有足够内射对象的 Abel 范畴，因此可以把空间 $X$ 上的层上同调群定义为如下函子的右导出函子：该函子把层映到其整体截面组成的 Abel 群.固定系数群 $\Lambda$ 的上同调，由它定义的常值层的上同调恢复.这现在是上同调群公认的定义，也是我们定义平展上同调群所采用的方法.我们不固定系数群并为了刻画上同调群而必须考虑所有(容许的)拓扑空间对，而是固定拓扑空间，但考虑该空间上的所有层.

\subsection*{层上同调简述}

设 $X$ 是拓扑空间.我们把 $X$ 的开子集做成一个范畴，其态射只有包含映射，并定义预层为从该范畴到 Abel 群范畴 $\mathbf{Ab}$ 的反变函子.于是，这样的预层 $\mathcal{F}$ 给 $X$ 的每个开子集 $U$ 配上一个 Abel 群 $\mathcal{F}(U)$，给每个包含 $V \subset U$ 配上限制映射 $\rho^{U}_{V} \colon \mathcal{F}(U) \to \mathcal{F}(V)$，使得 $\rho^{U}_{U} = \mathrm{id}_{\mathcal{F}(U)}$，并且每当 $W \subset V \subset U$ 时
\[
\rho^{U}_{W} = \rho^{V}_{W} \circ \rho^{U}_{V}.
\]
由于历史原因，$\mathcal{F}(U)$ 的元素称为 $\mathcal{F}$ 在 $U$ 上的截面，$\mathcal{F}(X)$ 的元素称为 $\mathcal{F}$ 的整体截面.此外，有时把 $\mathcal{F}(U)$ 写为 $\Gamma(U, \mathcal{F})$，把 $\rho^{U}_{V}(s)$ 写为 $s|_{V}$.

预层 $\mathcal{F}$ 称为层，如果
\begin{enumerate}
  \item 截面 $f \in \mathcal{F}(U)$ 由它在 $U$ 的开覆盖 $(U_i)_{i \in I}$ 的各集合上的限制 $\rho^{U}_{U_i}(f)$ 决定；
  \item 对 $U$ 的开覆盖 $(U_i)_{i \in I}$，一族截面 $f_i \in \mathcal{F}(U_i)$ 由某个截面 $f \in \mathcal{F}(U)$ 限制而来，当且仅当对所有 $i, j$，$f_i|_{U_i \cap U_j} = f_j|_{U_i \cap U_j}$.
\end{enumerate}
换言之，$\mathcal{F}$ 是层，如果对 $X$ 的开子集 $U$ 的每个开覆盖 $(U_i)_{i \in I}$，序列
\[
\mathcal{F}(U) \to \prod_{i \in I} \mathcal{F}(U_i) \rightrightarrows \prod_{(i,j) \in I \times I} \mathcal{F}(U_i \cap U_j)
\]
正合---按定义，这意味着第一个箭头把 $\mathcal{F}(U)$ 单射地映到 $\prod \mathcal{F}(U_i)$ 中在其上后两个箭头相等的子集上.第一个箭头把 $f \in \mathcal{F}(U)$ 送到族 $(f|_{U_i})_{i \in I}$，后两个箭头分别把 $(f_i)_{i \in I}$ 送到族 $(f_i|_{U_i \cap U_j})_{(i,j) \in I \times I}$ 与 $(f_j|_{U_i \cap U_j})_{(i,j) \in I \times I}$.由于我们只考虑 Abel 群的(预)层，可以把条件重述为：序列
\[
0 \to \mathcal{F}(U) \to \prod_{i \in I} \mathcal{F}(U_i) \to \prod_{(i,j) \in I \times I} \mathcal{F}(U_i \cap U_j), \qquad
f \mapsto (f|_{U_i}); \quad (f_i) \mapsto (f_j|_{U_i \cap U_j} - f_i|_{U_i \cap U_j})
\]
正合.把该条件应用于空集的空覆盖，蕴含 $\mathcal{F}(\varnothing) = 0$.\footnote{空集 $\varnothing$ 被一族以空集为指标集的映射覆盖，因此层条件说 $\mathcal{F}(\varnothing)$ 等于一个以 $\varnothing$ 为指标集的 Abel 群乘积.但在每个范畴中，空指标集上的乘积就是终对象(当它存在时)，对 Abel 群而言它为零.}

例如，若 $\Lambda$ 是拓扑 Abel 群(如 $\mathbb{R}$ 或 $\mathbb{C}$)，则可以在任意拓扑空间 $X$ 上定义层，方法是令 $\mathcal{F}(U)$ 等于连续映射 $U \to \Lambda$ 的集合，限制映射取函数的通常限制.

当 $\Lambda$ 带离散拓扑时，每个连续映射 $f \colon U \to \Lambda$ 在 $U$ 的每个连通分支上都是常值的，从而穿过 $\pi_0(U)$(即 $U$ 的连通分支空间)分解.当后一个空间离散时，$\mathcal{F}(U)$ 是所有映射 $\pi_0(U) \to \Lambda$ 的集合，即 $\mathcal{F}(U) = \Lambda^{\pi_0(U)}$.在这种情况下，我们称 $\mathcal{F}$ 是由 Abel 群 $\Lambda$ 定义的常值层.

Grothendieck 证明，在自然结构下，$X$ 上的层构成 Abel 范畴.于是，我们有内射层的概念：内射层是层 $\mathcal{I}$，使得对任意层 $\mathcal{F}$ 的任意子层 $\mathcal{F}_0$，每个同态 $\mathcal{F}_0 \to \mathcal{I}$ 都延拓为同态 $\mathcal{F} \to \mathcal{I}$.Grothendieck 证明，每个层都能嵌入内射层.

函子 $\mathcal{F} \mapsto \mathcal{F}(X)$(从 $X$ 上的层范畴到 Abel 群范畴)是左正合的，但一般不是右正合的.我们定义 $H^{r}(X, \Lambda)$ 为它的第 $r$ 个右导出函子.于是，给定层 $\mathcal{F}$，我们选择正合序列
\[
0 \to \mathcal{F} \to \mathcal{I}^{0} \to \mathcal{I}^{1} \to \mathcal{I}^{2} \to \cdots
\]
使每个 $\mathcal{I}^{r}$ 内射，并令 $H^{r}(X, \mathcal{F})$ 等于 Abel 群复形
\[
\mathcal{I}^{0}(X) \to \mathcal{I}^{1}(X) \to \mathcal{I}^{2}(X) \to \cdots
\]
的第 $r$ 个上同调群.

内射分解对定义上同调群是有用的，但对计算它不方便.作为替代，如果限制映射 $\mathcal{F}(U) \to \mathcal{F}(V)$ 对所有开集 $U \supset V$ 都是满射，我们称层 $\mathcal{F}$ 为松的(flabby)，并且证明若 $\mathcal{F}$ 松，则 $H^{r}(X, \mathcal{F}) = 0$.因此，用松层的分解可以计算上同调.

\subsection*{Zariski 拓扑的不足}

如我们上面所述，对许多目的而言，Zariski 拓扑的开子集太少.我们列出一些这种情况明显可见的情形.

\textbf{上同调群为零} 回忆拓扑空间 $X$ 称为不可约的，如果 $X$ 的任意两个非空开子集有非空交，而簇(或概形)称为不可约的，如果它作为(Zariski)拓扑空间不可约.

\begin{theorem}\label{thm:1.1}[Grothendieck 定理]
若 $X$ 是不可约拓扑空间，则对所有常值层和所有 $r > 0$，$H^{r}(X, \mathcal{F}) = 0$.
\end{theorem}

\begin{proof}
显然，不可约拓扑空间的每个开子集 $U$ 都是连通的.因此，若 $\mathcal{F}$ 是由群 $\Lambda$ 定义的常值层，则对每个非空开集 $U$，$\mathcal{F}(U) = \Lambda$.特别地，$\mathcal{F}$ 是松的，从而对 $r > 0$，$H^{r}(X, \mathcal{F}) = 0$.
\end{proof}

\begin{remark}\label{rem:1.2}
\v{C}ech 上同调群也为零.设 $\mathcal{U} = (U_i)_{i \in I}$ 是 $X$ 的开覆盖.则覆盖 $\mathcal{U}$ 的 \v{C}ech 上同调群是一个复形的上同调群，其第 $r$ 个群为
\[
\prod_{(i_0, \ldots, i_r) \in I^{r+1},\, U_{i_0} \cap \cdots \cap U_{i_r} \neq \varnothing} \Lambda,
\]
映射是显然的.对不可约空间 $X$，这个复形与空间 $X$ 无关；事实上，它只依赖于 $I$ 的基数(假定 $U_i$ 非空).容易构造该复形的收缩同伦，从而推出该复形正合.
\end{remark}

\textbf{逆映射定理不成立} $\mathcal{C}^{1}$ 映射 $\varphi \colon N \to M$(可微流形之间的)称为在 $n \in N$ 处平展，如果切空间上的映射 $d\varphi \colon \operatorname{Tgt}_{n}(N) \to \operatorname{Tgt}_{\varphi(n)}(M)$ 是同构.

\begin{theorem}\label{thm:1.3}[逆映射定理]
可微流形的 $\mathcal{C}^{1}$ 映射在其平展的每一点处都是局部同构，即若 $\varphi \colon N \to M$ 在 $n \in N$ 处平展，则存在 $n$ 与 $\varphi(n)$ 的开邻域 $V$ 与 $U$，使得 $\varphi$ 限制为同构 $V \to U$.
\end{theorem}

设 $X$ 与 $Y$ 是代数闭域 $k$ 上的非奇异代数簇.正则映射 $\varphi \colon Y \to X$ 称为在 $y \in Y$ 处平展，如果 $d\varphi \colon \operatorname{Tgt}_{y}(Y) \to \operatorname{Tgt}_{\varphi(y)}(X)$ 是同构.

例如，正如我们很快将看到的，$x \mapsto x^{n} \colon \mathbb{A}^{1}_{k} \to \mathbb{A}^{1}_{k}$ 除原点外处处平展(只要 $n$ 不被 $k$ 的特征整除).但是，若 $n > 1$，这个映射在任意一点都不局部同构；事实上，不存在 $\mathbb{A}^{1}_{k}$ 的非空开子集 $V$ 与 $U$，使得映射 $x \mapsto x^{n}$ 把 $V$ 同构地映到 $U$ 上.要看出这一点，注意 $x \mapsto x^{n}$ 对应于 $k$-代数同态 $T \mapsto T^{n} \colon k[T] \to k[T]$.若 $x \mapsto x^{n}$ 把 $V$ 映入 $U$，则函数域上的相应映射为 $T \mapsto T^{n} \colon k(T) \to k(T)$.若 $V \to U$ 是同构，则函数域上的映射也应是同构，但它不是.

取 $n = 2$，$k = \mathbb{C}$，于是映射是 $z \mapsto z^{2} \colon \mathbb{C} \to \mathbb{C}$.要证明它在 $z_0 \neq 0$ 处是局部同构，我们必须在包含 $z_0^{2}$ 的某个开集上构造 $z \mapsto z^{2}$ 的逆函数 $z \mapsto \sqrt{z}$.为了复解析地做到这一点，我们需要在复平面中从 0 到 $\infty$ 去掉一条曲线.这样的曲线的补集对复拓扑是开的，但对 Zariski 拓扑不是.

\textbf{纤维丛不局部平凡} 集合上的拓扑使我们能够谈论某件事在``点附近''而非``在点处''为真，以及在每一点的邻域中``局部地''为真.

例如，假设给定了代数闭域 $k$ 上的簇的正则映射 $\varphi \colon Y \to X$，并且每个纤维 $\varphi^{-1}(x)$ 上有 $k$-向量空间结构.$\varphi^{-1}(x)$ 的一组基的选择确定了一个 $k$-线性代数簇同构 $\varphi^{-1}(x) \to \mathbb{A}^{n}_{k}$($n$ 适当).映射 $\varphi \colon Y \to X$ 称为向量丛，如果它局部平凡，即每个点 $x \in X$ 都有开邻域 $U$，使得有交换图
\[
\begin{tikzcd}[row sep=small]
  \varphi^{-1}(U) \arrow[r] \arrow[rd] & U \times \mathbb{A}^{n} \arrow[d] \\
                                       & U
\end{tikzcd}
\]
其中顶上的箭头是 $k$-线性代数簇同构.

对 $\mathbb{C}$ 上的光滑簇，上述那种正则映射 $Y \to X$ 若对复拓扑局部平凡，则对 Zariski 拓扑局部平凡.因此 Zariski 拓扑对向量丛的研究是足够细的(见 Weil, A., \emph{Fibre spaces in algebraic geometry}, Notes of a course 1952, University of Chicago).然而，对更奇异的丛的研究，它不够细.考虑例如一个正则映射 $\pi \colon Y \to X$，其上赋予代数群 $G$ 在 $Y$ 上的作用(在 $X$ 上)，即满足群作用通常条件的正则映射 $m \colon Y \times G \to Y$，并且使得 $\pi(yg) = \pi(y)$.在 1958 年 S\'eminaire Chevalley 的一次报告中，Serre 把这样的系统称为局部等平凡的(locally isotrivial)，如果对每个 $x \in X$，存在有限平展映射 $U \to U_0 \subset X$，其中 $U_0$ 是 $x$ 的 Zariski 开邻域，使得 $Y$ 到 $U$ 的拉回同构于平凡系统 $U \times G \to U$.

这个概念的有用性，连同 Weil 的观察，引导 Grothendieck 引入平展拓扑.

\subsection*{平展上同调}

设 $X$ 与 $Y$ 是代数闭域 $k$ 上的光滑簇.正则映射 $\varphi \colon Y \to X$ 称为平展的，如果它在所有点 $y \in Y$ 处都平展.平展映射是拟有限的(其纤维是有限的)且开的.

$X$ 上的平展``拓扑''是这样的拓扑：其中``开集''是平展态射 $U \to X$.一族平展态射 $(U_i \to U)_{i \in I}$(在 $X$ 上)是 $U$ 的覆盖，如果 $U = \bigcup \varphi_i(U_i)$.点 $x \in X$ 的平展邻域是平展映射 $U \to X$ 连同 $U$ 中映到 $x$ 的点 $u$.

定义 $\mathbf{Et}/X$ 为如下范畴：对象是平展映射 $U \to X$，箭头是交换图
\[
\begin{tikzcd}[row sep=small]
  V \arrow[rd] \arrow[rr] & & U \arrow[ld] \\
                          & X &
\end{tikzcd}
\]
其中映射 $V \to X$ 与 $U \to X$ 都是平展的(于是 $V \to U$ 也自动平展).

$X$ 上平展拓扑的预层是反变函子 $\mathcal{F} \colon \mathbf{Et}/X \to \mathbf{Ab}$.如果对一切平展覆盖 $(U_i \to U)$，序列
\[
\mathcal{F}(U) \to \prod_{i \in I} \mathcal{F}(U_i) \rightrightarrows \prod_{(i,j) \in I \times I} \mathcal{F}(U_i \times_{U} U_j)
\]
正合，则它是层.与经典情形类似地证明，层范畴是 Abel 范畴，且有足够的内射对象.因此，可以完全像经典情形那样，用函子 $\mathcal{F} \mapsto \mathcal{F}(X)$ 的导出函子定义平展上同调群 $H^{r}(X_{\mathrm{et}}, \mathcal{F})$.

课程的其余部分将致力于研究这些群及其应用.

\subsection*{与复拓扑的比较}

设 $X$ 是 $\mathbb{C}$ 上的非奇异代数簇.$X$ 上的平展拓扑与 $X(\mathbb{C})$ 上的复拓扑哪个更细？严格地说，这个问题没有意义，因为它们不是同类型的对象，但复逆映射定理表明，点 $x \in X$ 的每个平展邻域都``包含''一个复邻域.更精确地说，给定 $x$ 的平展邻域 $(U, u) \to (X, x)$，存在 $x$ 对复拓扑的邻域 $(V, x)$，使得包含 $(V, x) \hookrightarrow (X, x)$ 分解为
\[
(V, x) \to (U, u) \to (X, x)
\]
(其中第一个映射是复解析的，第二个是代数的).

于是，$X$ 的每个平展覆盖都能被一个复拓扑的覆盖``加细''.由此，对任意 Abel 群 $\Lambda$，得到典范映射 $H^{r}(X_{\mathrm{et}}, \Lambda) \to H^{r}(X(\mathbb{C}), \Lambda)$.

\begin{theorem}\label{thm:1.4}[比较定理]
对每个有限 Abel 群 $\Lambda$，典范映射 $H^{r}(X_{\mathrm{et}}, \Lambda) \to H^{r}(X(\mathbb{C}), \Lambda)$ 都是同构.
\end{theorem}

取 $\Lambda = \mathbb{Z}/\ell^{n}\mathbb{Z}$ 并对 $n$ 取逆向极限，得到同构 $H^{r}(X_{\mathrm{et}}, \mathbb{Z}_{\ell}) \to H^{r}(X(\mathbb{C}), \mathbb{Z}_{\ell})$.与 $\mathbb{Q}_{\ell}$ 作张量积后，这变成同构 $H^{r}(X_{\mathrm{et}}, \mathbb{Q}_{\ell}) \to H^{r}(X(\mathbb{C}), \mathbb{Q}_{\ell})$.

\subsection*{平展上同调的应用}

平展上同调已成为算术几何，$\mathbb{C}$ 之外基域上的代数几何，数论的一部分，$K$ 理论的一部分，以及有限群与 $p$ 进群表示论的必备工具.我希望在课程中解释其中一些应用.这里我简要描述两个.

\textbf{经典代数几何} 给定域 $k$ 上的代数簇 $X$ 和域同态 $\sigma \colon k \to K$，得到 $K$ 上的代数簇 $^{\sigma}X$.例如，若 $X$ 是由方程
\[
\sum a_{i_1 \cdots i_n} T^{i_1}_{1} \cdots T^{i_n}_{n} = 0
\]
定义的仿射(或射影)簇，则 $^{\sigma}X$ 是由方程
\[
\sum \sigma a_{i_1 \cdots i_n} T^{i_1}_{1} \cdots T^{i_n}_{n} = 0
\]
定义的簇.注意我们还得到一个映射 $(t_1, \ldots, t_n) \mapsto (\sigma t_1, \ldots, \sigma t_n) \colon X(k) \to X(K)$.

\begin{question}\label{question:1.5}
设 $X$ 是 $\mathbb{C}$ 上的簇，$\sigma$ 是 $\mathbb{C}$ 的自同构(仅作为域，即不要求连续).上同调群 $H^{r}(X(\mathbb{C}), \mathbb{Q})$ 与 $H^{r}({}^{\sigma}X(\mathbb{C}), \mathbb{Q})$ 之间有什么关系？
\end{question}

映射 $(t_1, \ldots, t_n) \mapsto (\sigma t_1, \ldots, \sigma t_n) \colon X(\mathbb{C}) \to X(\mathbb{C})$ 不会是同胚，除非 $\sigma$ 是恒等映射或复共轭.事实上，Serre\footnote{\emph{Exemples de vari\'et\'es projectives conjugu\'ees non hom\'eomorphes}. C. R. Acad. Sci. Paris 258 (1964), 4194--4196.}构造了 $\mathbb{C}$ 上的非奇异射影曲面 $X$，使得 $X(\mathbb{C})$ 与 $({}^{\sigma}X)(\mathbb{C})$ 的基本群不同构，从而任何映射都不是同胚.Shimura 簇理论给出其他例子：簇 $X$ 使得 $X(\mathbb{C})$ 与 $({}^{\sigma}X)(\mathbb{C})$ 的基本群不同构.\footnote{见 Milne, James S.; Suh, Junecue. \emph{Nonhomeomorphic conjugates of connected Shimura varieties}. Amer. J. Math. 132 (2010), no. 3, 731--750. arXiv:0804.1953}

问题 \ref{question:1.5} 的答案是：群 $H^{r}(X(\mathbb{C}), \mathbb{Q})$ 与 $H^{r}({}^{\sigma}X(\mathbb{C}), \mathbb{Q})$ 在与 $\mathbb{Q}_{\ell}$(或 $\mathbb{C}$)作张量积后典范同构.

为证明这一点，使用比较定理以及 $H^{r}(X_{\mathrm{et}}, \mathbb{Q}_{\ell})$ 典范同构于 $H^{r}(({}^{\sigma}X)_{\mathrm{et}}, \mathbb{Q}_{\ell})$ 这一事实(因为同构 $\sigma \colon \mathbb{C} \to \mathbb{C}$ 把平展映射 $U \to X$ 送到平展映射 $^{\sigma}U \to {}^{\sigma}X$ 等---它是代数同构，而平展上同调群的定义是纯代数的).

这蕴含 $X$ 与 $^{\sigma}X$ 有相同的 Betti 数，从而 $H^{r}(X(\mathbb{C}), \mathbb{Q})$ 与 $H^{r}(({}^{\sigma}X)(\mathbb{C}), \mathbb{Q})$ 同构(但非典范同构).对光滑射影簇，这个较弱的命题由 Serre\footnote{\emph{G\'eom\'etrie alg\'ebrique et g\'eom\'etrie analytique}, Ann. Inst. Fourier, Grenoble 6 (1956), pp 1--42.}利用 Dolbeault 定理
\[
\beta_r(X) = \sum_{p+q=r} \dim H^{q}(X(\mathbb{C}), \Omega^{p})
\]
和他自己的比较定理
\[
H^{r}(X, \mathcal{F}) = H^{r}(X(\mathbb{C}), \mathcal{F}^{\mathrm{h}})
\]
证明(其中 $\mathcal{F}$ 是 $X$ 上对 Zariski 拓扑的 $\mathcal{O}_X$-模凝聚层，$\mathcal{F}^{\mathrm{h}}$ 是 $X(\mathbb{C})$ 上对复拓扑的 $\mathcal{O}^{\mathrm{h}}_{X}$-模的相伴层).一般簇的情形直到 Artin 证明比较定理之后才为人所知.

现在我可以解释为什么应当预期群 $H^{r}(X_{\mathrm{et}}, \mathbb{Q})$ 是异常的.给定 $\mathbb{C}$ 上的簇 $X$，存在 $\mathbb{C}$ 的子域 $k$，使得 $\mathbb{C}$ 是 $k$ 的无限 Galois 扩张，并且 $X$ 可以由系数在 $k$ 中的方程定义.于是对所有 $\sigma \in \operatorname{Gal}(\mathbb{C}/k)$, $^{\sigma}X = X$，从而 $\operatorname{Gal}(\mathbb{C}/k)$ 将作用在群 $H^{r}(X_{\mathrm{et}}, \mathbb{Q}_{\ell})$ 上.可以预期(事实上 Tate 猜想蕴含) $\operatorname{Gal}(\mathbb{C}/k)$ 在 $\operatorname{Aut}(H^{r}(X_{\mathrm{et}}, \mathbb{Q}_{\ell}))$ 中的像通常是无限的(从而是不可数的).由于 $H^{r}(X_{\mathrm{et}}, \mathbb{Q})$ 是 $\mathbb{Q}$ 上的有限维向量空间，$\operatorname{Aut}(H^{r}(X_{\mathrm{et}}, \mathbb{Q}))$ 是可数的，所以 $\operatorname{Gal}(\mathbb{C}/k)$ 在 $H^{r}(X_{\mathrm{et}}, \mathbb{Q}_{\ell})$ 上的作用不可能来自 $H^{r}(X_{\mathrm{et}}, \mathbb{Q})$ 上的作用.

\textbf{有限群的表示论} 设 $G$ 是有限群，$k$ 是特征零的域.群代数 $k[G]$ 是半单 $k$-代数，因此存在有限集 $\{V_1, \ldots, V_r\}$ 的简单 $k[G]$-模，使得每个 $k[G]$-模都同构于若干 $V_i$ 的直和，且每个简单 $k[G]$-模恰好同构于其中一个 $V_i$.事实上，人们知道 $r$ 等于 $G$ 中共轭类的个数.

对每个 $i$，定义 $\chi_i(g)$($g \in G$)为 $g$ 作用在 $V_i$ 上的迹.则 $\chi_i(g)$ 只依赖于 $g$ 的共轭类.$G$ 的特征标表列出了 $\chi_1, \ldots, \chi_r$ 在每个共轭类上的值.人们希望计算每个有限单群的特征标表.

对交换单群，这是平凡的.$A_n$ 的表示论与 $S_n$ 的密切相关，由 Frobenius 和 Young 在 1900 到 1930 年间完成.散在群的特征标表是已知的---毕竟它们只有 26 个.剩下的就是与代数群相关的那些(迄今为止最大的一类).

$\operatorname{GL}_n(\mathbb{F}_q)$ 的表示论由 J.A. Green 完成，$\operatorname{Sp}_4(\mathbb{F}_q)$ 的由 B. Srinivasan 完成.

设 $G$ 是 $\mathbb{F}_q$ 上的代数群.则 $G(\mathbb{F}_q)$ 是 Frobenius 映射 $t \mapsto t^{q}$ 作用在 $G(\mathbb{F}_q^{\mathrm{al}})$ 上的不动点集：
\[
G(\mathbb{F}_q) = G(\mathbb{F}_q^{\mathrm{al}})^{F}.
\]
为了得到所有有限单群，需要考虑 $\mathbb{F}_p^{\mathrm{al}}$ 上的约化群 $G$ 和映射 $F \colon G(\mathbb{F}_p^{\mathrm{al}}) \to G(\mathbb{F}_p^{\mathrm{al}})$，使得 $F$ 的某个幂 $F^{m}$ 是 $G$ 在 $\overline{\mathbb{F}}_q$ 上某个模型的 Frobenius 映射($q$ 适当)；有限群再次是 $F$ 在 $G(\mathbb{F}_p^{\mathrm{al}})$ 上的不动点集.

Deligne 和 Lusztig 1976 以及 Lusztig 之后的大量论文，非常成功地把平展上同调应用于研究这些群的表示论.他们构造了 $\mathbb{F}$ 上 $G$ 作用的簇 $X$.于是 $G$ 作用于 $X$ 的平展上同调群，并且可以用 Lefschetz 不动点公式计算这些表示的迹.
